\section{Findings \& Discussion of Results}
\label{sec:Fin}

This section reports the experimental outcomes of the server-side ECU trace replay with LARVA instrumentation toggled on/off. Results are drawn from the benchmark session with tuned \texttt{loadMapMismatch} threshold (\texttt{|Load-MAPSource| > 60}), using 3 warmup runs and 10 measured runs per mode on the same machine and JVM conditions. The experiment assumes telemetry is transmitted from an IoT/ECU endpoint to a server-side verification API.

\subsection{Result presentation}

Two regimes were evaluated:
\begin{itemize}
    \item \textbf{Baseline regime:} clean trace (\texttt{20180713-home2mimos\_clean.csv}).
    \item \textbf{Abnormal regime:} injected trace (\texttt{20180713-home2mimos\_clean\_abnormal.csv}) with approximately 20\% violation rows across four flags.
\end{itemize}

Table~\ref{tab:rv-overhead} summarises core timing and throughput metrics. Standard deviation is reported with each average to capture run-to-run variability.

\begin{table}[ht]
    \centering
    \caption{Runtime verification overhead summary (10 measured runs per mode)}
    \label{tab:rv-overhead}
    \begin{tabular}{l|c|c}
    \textbf{Metric} & \textbf{Baseline} & \textbf{Abnormal} \\ \hline
    RV ON elapsed (ms) & 113.0 $\pm$ 7.18 & 4274.8 $\pm$ 137.40 \\
    RV OFF elapsed (ms) & 104.3 $\pm$ 8.39 & 105.5 $\pm$ 4.09 \\
    RV ON throughput (rows/s) & 607{,}526.6 $\pm$ 33{,}682.8 & 16{,}023.1 $\pm$ 506.6 \\
    RV OFF throughput (rows/s) & 659{,}328.3 $\pm$ 44{,}358.5 & 649{,}524.9 $\pm$ 24{,}972.4 \\
    Overhead (\%) & 8.34 & 3951.94 \\
    \end{tabular}
\end{table}

Tail latency also reflects this split behaviour. For baseline, RV ON p95/p99 are 124.0/131.2 ms, while RV OFF p95/p99 are 117.2/125.84 ms. For abnormal traces, RV ON p95/p99 rise to 4497.6/4530.72 ms, whereas RV OFF remains at 111.0/111.0 ms.

\subsection{Analysis and interpretation}

\textbf{RQ1 (detection of injected abnormalities).} The abnormal regime produced high RV ON flag totals (\texttt{coolantTooHigh}=34210, \texttt{voltageOutOfRange}=34200, \texttt{throttleSpike}=34220, \texttt{loadMapMismatch}=34230), while RV OFF flags remained zero. This indicates that violation reporting is coupled to monitor execution as intended and that injected abnormal classes are detected consistently.

\textbf{RQ2 (monitoring overhead).} Baseline overhead is moderate (8.34\%), suggesting that monitor integration is feasible for non-adversarial trace processing on the tested server path. However, abnormal overhead is very high (3951.94\%), showing that violation-heavy execution paths can dominate runtime. In other words, monitor cost is not constant; it is strongly state-dependent. These overhead values should be interpreted as server-side verification cost under a telemetry-upload architecture.

\textbf{RQ3 (stability and threshold sensitivity).} Standard deviations are small relative to means in most cases, indicating repeatable behaviour. After tuning \texttt{loadMapMismatch} from 30 to 60, baseline false positives reduced (aggregate baseline \texttt{loadMapMismatch} decreased compared with pre-tuning runs), while abnormal detection remained strong. This supports threshold calibration as a practical control for false-positive pressure.

\subsection{Comparison with related work and quality considerations}

The observed trade-off aligns with prior RV literature: runtime monitoring provides actionable safety/cybersecurity visibility, but overhead depends on property design, trigger frequency, and deployment context~\cite{krichen2023survey,zhang2024context,jaksic2018algebraic}. The clear difference between baseline and abnormal runtime also mirrors broader automotive V\&V concerns, where worst-case and abuse-case behaviour must be measured rather than assumed from nominal conditions~\cite{ekert2021cybersecurity}.

\textbf{Validity.} Internal validity is strengthened by same-endpoint comparisons, warmup before measured runs, and on/off toggling as the only intended treatment factor. Construct validity is supported by explicit metrics (elapsed time, throughput, p95/p99, overhead) and deterministic replay inputs.

\textbf{Reliability.} Repeated 10-run measurements with reported standard deviations indicate reproducibility under stable local conditions.

\textbf{Generalizability.} External validity remains limited to this server-side replay implementation and selected property set. Results should not be extrapolated directly to production ECUs without hardware-in-the-loop or in-vehicle evaluation.